\documentclass[10pt,a4paper]{amsart}
\usepackage[margin=19mm]{geometry}
\usepackage{amsmath,amssymb,mathtools}
\usepackage[scr=rsfs,cal=euler]{mathalfa}
\usepackage{xfrac}
\usepackage[unicode,hidelinks,bookmarks=false]{hyperref}
\newcommand{\ie}{i.e.\ }
\newcommand{\cf}{cf.\ }
\newcommand{\resp}{respectively\ }
\newcommand{\Subsection}{\subsection}
\providecommand{\nobreakdash}{\nobreak\hbox{-}}
\setlength{\emergencystretch}{2em}
\multlinegap0pt
\usepackage{mathtools}
\usepackage{amssymb}
\usepackage{booktabs}
\newtheorem{proposition}{Proposition}
\title{Paley-Wiener Theorem for Probabilistic Frames}
\author{Dongwei Chen}
\date{Source: arXiv:2310.17830v3 (CC BY 4.0)}
\begin{document}
\maketitle

\noindent\emph{Source and licence.} Paley-Wiener Theorem for Probabilistic Frames. Version arXiv:2310.17830v3 (CC BY 4.0), distributed by arXiv under the Creative Commons Attribution 4.0 International licence, http://creativecommons.org/licenses/by/4.0/, as recorded in the arXiv licence metadata for this exact version and shown on the abstract page https://arxiv.org/abs/2310.17830v3. Full licence notice: \texttt{licenses/arxiv-2310.17830-CC-BY-4.0.txt}.\par\smallskip

\noindent\textit{Editorial scope.} Counted unit: the article's proposition CoulingDualPerturbation (source0.tex lines 509--520). The article's complete original proof of it is delivered with it (source0.tex lines 522--565, one \texttt{proof} environment of 44 lines). No mathematics is written, repaired or re-derived: the statement and the proof are the article's own lines, in the article's own notation, and no retained line is substituted or re-flowed.  Retained uncounted as the source-local context the proof needs: lines 206--208.  Nothing was omitted from any retained span, so the package carries no omission record.  No retained line contains a citation, so the wrapper carries no bibliography and no bibliography entry.  No retained line contains a figure environment, an image-inclusion command or drawing code, so no image is delivered and the package declares no asset dependency.  Editorial formatting only, in addition: an AMS \texttt{amsart} preamble that restates the article's own declarations and macros used by the retained lines, namely the environments \texttt{proposition} on separate counters, together with the standard packages those lines need.  The retained environments print in this document as Proposition 1 in order of appearance, whereas the article prints the same items with the numbering of its own full document.  The complete native source of the article as distributed in the arXiv e-print for this exact version is bundled unchanged as \texttt{sources/149142/source0.tex}.
\begin{equation}\label{reconstruction2}
   {\bf f}=   \int_{\mathbb{R}^n} \langle  {\bf f}, {\bf S}_\mu^{-1} {\bf x} \rangle {\bf x}  d\mu({\bf x}) =\int_{\mathbb{R}^n} \langle  {\bf S}_\mu^{-1} {\bf f}, {\bf x} \rangle {\bf x}  d\mu({\bf x}). 
\end{equation}

\begin{proposition}\label{CoulingDualPerturbation}
Let $\mu$ be a probabilistic frame for $\mathbb{R}^n$ with bounds $0<A \leq B$ and $\eta \in \mathcal{P}_2(\mathbb{R}^n)$. Suppose there exists $\gamma \in \Gamma(\mu, \eta)$ such that
\begin{equation*}
   \epsilon:= \int_{\mathbb{R}^n \times \mathbb{R}^n}  \Vert {\bf x}\Vert  \Vert {\bf x}-{\bf z} \Vert  d \gamma({\bf x},{\bf z})<A,
\end{equation*}
then $\eta$ is a probabilistic frame for $\mathbb{R}^n$ with bounds $\frac{(A- \epsilon)^2}{B} \ \text{and} \ M_2(\eta)$,
and if 
\begin{equation*}
   \chi := \int_{\mathbb{R}^n \times \mathbb{R}^n} \Vert {\bf S}_\mu^{-1} {\bf x}\Vert   \Vert {\bf x}-{\bf z} \Vert  d \gamma({\bf x},{\bf z})<1,
\end{equation*}
then $\eta$ is a probabilistic frame for $\mathbb{R}^n$ with bounds $\frac{A^2 (1-\chi)^2}{B} \ \text{and} \ M_2(\eta)$.
\end{proposition}

\begin{proof}
Since $\eta \in \mathcal{P}_2(\mathbb{R}^n)$, $\eta$ is Bessel with bound  $M_2(\eta)$.  Using the reconstruction formula for the canonical dual frame in Equation \eqref{reconstruction2} and $\gamma \in \Gamma(\mu, \eta)$, we have
$$ 
{\bf f} =\int_{\mathbb{R}^n} \langle {\bf S}_\mu^{-1} {\bf f}, {\bf x} \rangle {\bf x}  d\mu({\bf x}) = \int_{\mathbb{R}^n \times \mathbb{R}^n} \langle {\bf S}_\mu^{-1} {\bf f}, {\bf x} \rangle {\bf x} d \gamma({\bf x},{\bf z})  , \ \text{for any} \ {\bf f} \in \mathbb{R}^n.
$$
Now define $L: \mathbb{R}^n \rightarrow \mathbb{R}^n$ by 
     \begin{equation*}
         L({\bf f}) = \int_{\mathbb{R}^n \times \mathbb{R}^n} \langle {\bf S}_\mu^{-1} {\bf f}, {\bf x} \rangle {\bf z}  d \gamma({\bf x},{\bf z}), \ \text{for any} \ {\bf f} \in \mathbb{R}^n.
     \end{equation*}
Therefore, 
\begin{equation*}
    \Vert {\bf f} -L({\bf f}) \Vert 
    = \Big \Vert \int_{\mathbb{R}^n \times \mathbb{R}^n} \langle {\bf S}_\mu^{-1}{\bf f}, {\bf x} \rangle ({\bf x}-{\bf z}) d \gamma({\bf x}, {\bf z}) \Big \Vert 
     \leq \epsilon \ \Vert {\bf S}_\mu^{-1} \Vert_2 \ \Vert  {\bf f} \Vert  \leq \frac{\epsilon}{A}  \Vert  {\bf f} \Vert < \Vert  {\bf f} \Vert,
\end{equation*}
where the second inequality follows from that $\|{\bf S}_\mu^{-1}\|_2 \leq \frac{1}{A}$. 
Thus, $L: \mathbb{R}^n \rightarrow \mathbb{R}^n$ is invertible and $\Vert L^{-1} \Vert \leq \frac{1}{1-\epsilon/ A}$.
Then, for any ${\bf f} \in \mathbb{R}^n$, 
\begin{equation*}
    {\bf f} = LL^{-1}({\bf f}) =  \int_{\mathbb{R}^n \times \mathbb{R}^n} \langle {\bf S}_\mu^{-1}L^{-1} {\bf f}, {\bf x} \rangle {\bf z} d\gamma({\bf x},{\bf z}).
\end{equation*}
Therefore, 
\begin{equation*}
\begin{split}
   \Vert {\bf f} \Vert^4 &=  |\langle {\bf f}, {\bf f} \rangle|^2 = \Big \vert \int_{\mathbb{R}^n \times \mathbb{R}^n} \langle {\bf S}_\mu^{-1}L^{-1} {\bf f}, {\bf x} \rangle \langle  {\bf f}, {\bf z} \rangle d \gamma({\bf x},{\bf z}) \Big \vert^2 \\
    & \leq \int_{\mathbb{R}^n} \left|\langle {\bf S}_\mu^{-1} L^{-1}{\bf f}, {\bf x} \rangle \right|^2 d\mu({\bf x})  \ \int_{\mathbb{R}^n} \left| \langle {\bf f}, {\bf z} \rangle \right|^2 d\eta({\bf z})\\
 &\leq B \Vert  {\bf S}_\mu^{-1} \Vert_2^2  \ \Vert L^{-1} \Vert^2  \ \Vert {\bf f} \Vert^2   \int_{\mathbb{R}^n} \left|\langle {\bf f}, {\bf z} \rangle \right|^2 d\eta({\bf z}) \leq \frac{B\Vert {\bf f} \Vert^2}{A^2 (1-\frac{\epsilon}{A})^2 }  \int_{\mathbb{R}^n} \left| \langle {\bf f}, {\bf z} \rangle \right |^2 d\eta({\bf z}),
\end{split}
\end{equation*}
where the first inequality is due to the Cauchy-Schwarz inequality in $L^2(\gamma)$ and $\gamma \in \Gamma(\mu, \eta)$, and the second inequality follows from that $\mu$ is a probabilistic frame with upper bound $B$. 
Then, for any ${\bf f} \in \mathbb{R}^n$, 
$$ \frac{(A- \epsilon)^2}{B}    \Vert  {\bf f} \Vert^2  \leq   \int_{\mathbb{R}^n} |\langle {\bf f}, {\bf z} \rangle|^2 d\eta({\bf z}) \leq M_2(\eta) \Vert  {\bf f} \Vert^2.$$
Thus, $\eta$ is a probabilistic frame with bounds $\frac{(A- \epsilon)^2}{B}$ and $M_2(\eta)  $.
Furthermore, if 
\begin{equation*}
   \chi := \int_{\mathbb{R}^n \times \mathbb{R}^n}  \Vert {\bf S}_\mu^{-1} {\bf x}\Vert \Vert {\bf x}-{\bf z} \Vert  d \gamma({\bf x},{\bf z})<1,
\end{equation*}
then 
 \begin{equation*}
    \Vert {\bf f} -L({\bf f}) \Vert \leq  \int_{\mathbb{R}^n \times \mathbb{R}^n} |\langle {\bf f}, {\bf S}_\mu^{-1} {\bf x} \rangle| \ \|{\bf x}-{\bf z}\|  d \gamma({\bf x}, {\bf z}) \leq \chi \ \Vert  {\bf f} \Vert < \Vert  {\bf f} \Vert. 
\end{equation*}
Therefore, $L$ is invertible and $\Vert L^{-1} \Vert \leq \frac{1}{1-\chi}$.
Similar arguments show that $\eta$ is a probabilistic frame for $\mathbb{R}^n$ with bounds $\frac{A^2 (1-\chi)^2}{B} \ \text{and} \ M_2(\eta)$.
 \end{proof}

\end{document}
